\documentclass[12pt]{article}

\usepackage[utf8]{inputenc}
\usepackage[T1]{fontenc}

\title{Systematics}
\author{Vladimir Roa}
\author{Luis Gonzalez}
\date{\today}

\begin{document}

\maketitle

\section{Introducción}

La familia Labridae es la quinta más numerosa entre todos los peces, albergando alrededor de 600 especies distribuidas en cerca de 82 géneros (Westneat y Alfredo, 2005). Habitan en aguas costeras y de plataforma continetal en mares templados y tropicales, con un rango de tamaño que va desde unos 20 centimetros hasta ejemplares que superan los 100 kilogramos (Barber y Bellwood, 2005; Westneat y Alfredo, 2005). La diversidad anatomica de estos les permite tener una amplia alimentación y beneficiarse de muchos de los ecosistemas marinos (Westneat y Alfredo, 2005).Su clasificación clasica se ha dividido a los labridos en seis tribus: Hypsigenyini, Labrini, Cheilinini, Novaculini, Labrichthyini y Julidini (Westneat y Alfredo, 2005).

El propósito de este trabajo es evaluar la arquitectura filogenética de Labridae y profundizar en la sistemática de la tribu Julidini, con especial atención a la polifilia del genero *Halichoeres*. Mediante la integración de la hipotesis cladística moleculares, datación de divergencia y evidencia citogenética comparada, se analizan las relaciones evolutivas entre los componentes del Nuevo Mundo (Atlantico tropical y Pacifico oriental) y del viejo mundo (indopacífico y Atlantico oriental). También se constrastan los patrones de especiación y cladogenesis con eventos vicariantes de escala geológica, como el cierre de antiguio Mar de Tetis y la emergencia del Istmo de panamá.

Los analisis filogeneticos moleculares redefinieron los limites del grupo al demostrar que familias consideradas anteriormente independientes, como los Scaridae y los Odacidae, resultaron estar anidadas dentro de Labridae . Los odácidos forman parte de la tribu basal Hypsigenyini junto al género monotípico *Pseudodax moluccanus*, un pez asociado en el pasado a los peces loro por la fusión de sus dientes en forma de pico pero geneticamente cercano a los Lachnolaimus y Odacidae. Adicionalmente, se ha sugerido la creación de la tribu Pseudiolabrini para reunir a los géneros templados del pacífico como *Austrolabrus*, *Pictilabrus*, *Notolabrus* y *Pseudolabrus*(westneat y alfredo, 2005).


Referencias.
Barber, P. H., &amp; Bellwood, D. R. (2005). Biodiversity hotspots: Evolutionary origins of biodiversity in wrasses (Halichoeres: Labridae) in the Indo-Pacific and new world tropics. Molecular Phylogenetics and Evolution, 35(2), 235–253\. https://doi.org/10.1016/j.ympev.2004.10.004
Westneat, M. W., &amp; Alfaro, M. E. (2005). Phylogenetic relationships and evolutionary history of the reef fish family Labridae. Molecular Phylogenetics and Evolution, 36(2), 370–390\. https://doi.org/10.1016/j.ympev.2005.02.001

\end{document}
